\section[Problem]{The design problem}
\sectionframe{01}{The design problem}{Why ask a generative model for new microstructures?}{\faBalanceScale}

% ------------------------------------------------------------ the trade-off
\begin{frame}[source={Liao et al.\ (2026), Nat.\ Commun.\ 17, 11232}]{Strength/ductility trade-off}
  \begin{tikzpicture}[x=1mm,y=1mm]
    \useasboundingbox (0,0) rectangle (144,-54);
    \begin{scope}[shift={(16,-42)}]
      \planeaxes{340}{1000}{28}{86}{124}{42}{400,500,...,1000}{30,40,...,80}
      \begin{scope}[plane]
        \drawfront[front]
        \drawalloys
        \draw[udsred,line width=1.1pt,-{Stealth[length=2.6mm,width=2mm]}] (640,55) -- (760,72);
        \node[anchor=west,inner sep=0pt,font=\footnotesize\bfseries,text=udsred] at (768,73.5)
          {goal: stronger and more ductile};
        % the two ends of the trade-off, each beside the axis it is close to
        \node[anchor=west,inner sep=0pt,font=\footnotesize\bfseries,text=udsdark] at (410,84) {ductile, but weak};
        \node[anchor=east,inner sep=0pt,font=\footnotesize\bfseries,text=udsdark] at (1000,41) {strong, but brittle};
        \node[anchor=west,label text] at (360,33) {25 alloys, one processing route each};
        \node[anchor=west,small label] (fl) at (600,40) {measured front};
        \draw[udsgreymid,line width=0.4pt] (fl.north east) -- (700.46,45);
      \end{scope}
    \end{scope}
    \node[axis label] at (78,-49.4) {yield strength (MPa)};
    \node[axis label,rotate=90] at (3,-21) {elongation (\%)};
  \end{tikzpicture}
  \note{\textbf{1:00--2:10}\par
    Each point is one alloy: Inconel 625, a nickel-based superalloy used for turbine discs, blades and
      combustors, processed in one of 25 ways.\par
    The alloys differ only in processing: hot rolling and annealing at different temperatures and
      times.\par
    Each was tensile-tested for its yield strength and its elongation to failure.\par
    The stronger variants stretch less: a strength--ductility trade-off; making a metal stronger
      usually makes it brittle, and making it ductile makes it weak.\par
    The dashed staircase is the front: no measured alloy beats a point on it in both properties.\par
    The goal is a microstructure above and to the right of the front.\par
    Such microstructures are usually found by experience and trial and error over processing
      routes; the paper asks a generative model to propose them.}
\end{frame}

% ------------------------------------------------------- failure case
\begin{frame}[source={Liao et al.\ (2026), Nat.\ Commun.\ 17, 11232, Fig.\ 2c (CC BY-NC-ND 4.0)}]
  {Classic pixelwise loss fails}
  \begin{tikzpicture}[x=1mm,y=1mm]
    \useasboundingbox (0,0) rectangle (144,-54);
    \figpart{32}{0}{2}{340}{770}{1897}{1268}{112}
    \node[anchor=east,label text,align=right] at (29,{\fpy{884}}) {real patches};
    \node[anchor=east,label text,align=right,text=udsred] (pl) at (29,{\fpy{1155}}) {pixel loss only};
    % the two classic pixel losses over the n pixels of a patch x and its reconstruction
    \node[anchor=north east,inner sep=0pt,font=\scriptsize,text=udsdark,align=right] at ([yshift=-1.6mm]pl.south east)
      {MAE $=\frac{1}{n}\sum_i|x_i-\hat{x}_i|$\\[1.2mm] MSE $=\frac{1}{n}\sum_i(x_i-\hat{x}_i)^2$};
    \warnstrip{-41.4}{blurred or missing grain boundaries, colours mixed inside grains}
  \end{tikzpicture}
  \note{\textbf{2:10--3:10}\par
    These are real patches of 50 by 50 micrometres from test alloys, cut from EBSD maps.\par
    In an EBSD map each pixel's colour encodes the crystal orientation there: the three Euler angles
      as red, green and blue; the black lines are grain boundaries.\par
    The second row is what the paper's generator returns when it is trained to reproduce its input
      with a pixel loss alone, the mean absolute difference of the pixel values (MAE); the mean squared
      difference (MSE) is the other classic choice, and both average over all pixels.\par
    The boundaries blur, and colours mix inside grains, although a real grain has nearly one
      orientation.\par
    Boundaries cover few pixels, so an average over all pixels barely notices them; yet boundaries
      and orientations are what set the mechanical properties.}
\end{frame}

% ------------------------------------------------------------- the idea
\begin{frame}[source={Liao et al.\ (2026), Nat.\ Commun.\ 17, 11232, after Figs.\ 1 and 4f}]{Two fixes}
  \begin{tikzpicture}[x=1mm,y=1mm]
    \useasboundingbox (0,0) rectangle (144,-54);
    \tikzset{verdict mark/.style={circle,text=white,inner sep=0pt,minimum size=3.8mm,font=\scriptsize},
      head/.style={anchor=west,inner sep=0pt,font=\footnotesize\bfseries},
      panel title/.style={anchor=north west,inner sep=0pt,font=\footnotesize\bfseries,text=udsdark},
      note/.style={anchor=north west,label text,text width=22mm}}
    % column heads: what goes wrong, and what the paper changes
    \node[verdict mark,fill=udsred] (bad) at (1.9,-1.9) {\faTimes};
    \node[head,text=udsred] at ([xshift=1.6mm]bad.east) {the obstacle};
    \node[verdict mark,fill=testgreen] (good) at (87.9,-1.9) {\faCheck};
    \node[head,text=testgreen] at ([xshift=1.6mm]good.east) {the paper's fix};
    % row 1, the loss: the same patch region, three times enlarged, before and after
    \fill[redtint] (0,-5.6) rectangle (58,-28.4);
    \fill[greentint] (86,-5.6) rectangle (144,-28.4);
    \node[panel title] at (2.4,-7.4) {$L=L_{\mathrm{pix}}$};
    \node[panel title] at (88.4,-7.4) {$L=L_{\mathrm{pix}}+\textcolor{testgreen}{L_{\mathrm{edge}}+L_{\mathrm{SSIM}}}$};
    \foreach \x/\img in {2.4/grain-map-blurred,88.4/grain-map} {
      \node[anchor=north west,inner sep=0pt] at (\x,-12.6) {\includegraphics[width=14mm]{assets/drawn/\img}};
      \draw[udsdark,line width=0.4pt] ({\x+4.9},-17.5) rectangle ++(4.2,-4.2);
      \begin{scope}
        \clip ({\x+24},-19.6) circle (6.3);
        \node[anchor=north west,inner sep=0pt] at ({\x+3},1.4) {\includegraphics[width=42mm]{assets/drawn/\img}};
      \end{scope}
      \draw[udsdark,line width=0.5pt] ({\x+24},-19.6) circle (6.3);
      \draw[udsdark,line width=0.3pt] ({\x+9.1},-17.5) -- ({\x+19.6},-15.1) ({\x+9.1},-21.7) -- ({\x+19.6},-24.1);}
    \node[note] at (121.4,-13.4) {sharp boundaries, one colour per grain};
    \fixarrow{60.4}{83.6}{-17}{more terms}
    % row 2, the sampling: the latent space sampled at random, and searched toward targets (after Fig. 1c)
    \fill[redtint] (0,-31.2) rectangle (58,-54);
    \fill[greentint] (86,-31.2) rectangle (144,-54);
    \node[panel title] at (2.4,-33) {random sampling};
    \node[panel title] at (88.4,-33) {property-guided search};
    \tikzset{code/.style={circle,fill=udsgreymid,inner sep=0pt,minimum size=1.3mm},
      new code/.style={circle,fill=udsblue,inner sep=0pt,minimum size=1.5mm},
      latent/.style={fill=white,draw=udsgreylight,line width=0.4pt}}
    % left: codes drawn at random fall among the training codes
    \filldraw[latent] (14,-43.6) ellipse (10 and 7);
    \foreach \x/\y in {9.6/-38.6,7/-41.4,5.8/-43.8,8.4/-45,9.8/-48.4,14.6/-38.4,19.6/-39,18.6/-48.2,14.8/-49.2}
      \node[code] at (\x,\y) {};
    \node[anchor=north,small label] at (14,-51.2) {latent space};
    % the three new codes join, and one arrow goes on to the sentence about them; the training codes
    % leave each line a clear path
    \node[note,text width=27mm,anchor=west] at (28.4,-43.6) {new codes fall among the training codes};
    \foreach \x/\y in {12/-40.4,15.8/-43.6,11.2/-47}
      \draw[udsblue,line width=0.5pt,shorten <=1mm] (\x,\y) to[out=0,in=180] (22.4,-43.6);
    \draw[udsblue,line width=0.5pt,-{Stealth[length=1.5mm,width=1.1mm]}] (22.4,-43.6) -- (27.4,-43.6);
    \foreach \x/\y in {12/-40.4,15.8/-43.6,11.2/-47} \node[new code] at (\x,\y) {};
    % right: the search steers codes to targets outside them, beyond the measured front (after Fig. 1c)
    \filldraw[latent] (97.6,-45) ellipse (8 and 5.8);
    \foreach \x/\y in {-4.6/1.4,-2/-2.4,0.4/2.6,2.2/-0.6,4.6/1.6,-0.6/-4,5/-2.6,-5.4/-1.2}
      \node[code] at (97.6+\x,-45+\y) {};
    \draw[udsblue,line width=0.9pt,-{Stealth[length=1.8mm,width=1.4mm]}] (104.8,-42.6) -- (107.9,-41.4);
    \filldraw[fill=udsblue!22,draw=udsblue,line width=0.4pt] (111.2,-40.6) circle (3.2);
    \foreach \x/\y in {-1.2/0.9,1.2/-0.9} \node[text=udsred,font=\tiny] at (111.2+\x,-40.6+\y) {\faStar};
    \node[anchor=north,small label] at (97.6,-51.2) {latent space};
    \node[anchor=north,small label] at (111.2,-44.2) {targets};
    \begin{scope}[shift={(119,-49.4)}]
      \draw[udsgreymid,line width=0.5pt] (0,11.4) -- (0,0) -- (22,0);
      \node[anchor=north,small label] at (11,-0.6) {strength};
      \node[anchor=south,small label,rotate=90] at (-0.6,5.7) {elongation};
      \begin{scope}[x=0.033333mm,y=0.203571mm,shift={(-340,-30)}]
        \drawfront[front]
        \drawalloys[minimum size=1mm]
        \draw[udsred,line width=0.7pt,-{Stealth[length=1.6mm,width=1.2mm]}] (600,51) -- (660,61);
        \drawcandidates[minimum size=1.3mm]
      \end{scope}
    \end{scope}
    \fixarrow{60.4}{83.6}{-42.2}{steer by properties}
  \end{tikzpicture}
  \note{\textbf{3:10--4:10}\par
    The paper answers two obstacles with two changes: the obstacles on the left, its fixes on the
      right.\par
    First, a pixel loss alone blurs grain boundaries, as the enlarged region shows; the paper adds two
      terms, one that responds to grain boundaries and one to the structure inside grains.\par
    Second, codes drawn at random from the model's latent space fall among the training codes, so the new
      samples resemble the measured alloys.\par
    So the authors search the latent space, steered by predicted strength and ductility, toward targets
      outside the training codes, as their Fig.~1c sketches; their five candidates lie beyond the measured
      front.\par
    Section 2 explains the models and the search for anyone who has not met them.}
\end{frame}
